% =====================================================================
%  tr10a.tex —— 原文 pp.127--133（tr-06 追加块）
%  §22 附录：代数拓扑（同伦群、CW 复形、正合列等）、泛函分析、同调代数
%  说明：原文附录未使用任何编号定理类，全部为粗体小标题（Van Kampen's Theorem.
%        5-Lemma. 等），故按 STYLE_GUIDE §4 用 \para{…}，不自造编号。
% =====================================================================
% ===== 原文 pp.127--133 =====
% [原文 p.127]
\section{附录}

在本附录中，我们简要回顾来自数学其他分支的一些定义与结果，以帮助读者研习
（更多细节可参看引言中那份带背景材料的书目）。

我们从\emph{Seifert--van Kampen 定理}开始。该定理把某个拓扑空间（比如 $X$）的基本群，
用两个覆盖 $X$ 的开且道路连通的拓扑子空间 $X_1, X_2 \subset X$ 的基本群来表达
（因此它可用来计算那些由更简单的空间拼成的空间的基本群）：

\para{Van Kampen 定理.} 设 $X$ 是拓扑空间，它是两个道路连通子空间 $X_1, X_2 \subset X$ 的内部之并。
设 $X_0 = X_1 \cap X_2$ 非空且道路连通。又设 $*$ 是 $X_0$ 的一个基点，并设
$i_k : \pi_1(X_0, *) \to \pi_1(X_k, *)$ 与 $j_k : \pi_1(X_k, *) \to \pi_1(X, *)$
（$k = 1, 2$）是相应包含映射诱导的映射。则 $X$ 道路连通，且自然映射
\[
  \pi_1(X_1, *) \star_{\pi_1(X_0,*)} \pi_1(X_2, *) \to \pi_1(X, *)
\]
是同构，即 $X$ 的基本群是子空间 $X_1$ 与 $X_2$ 的基本群的自由积，并合并了 $X_0$ 的基本群。

证明的关键想法是：$X$ 中的道路可以分段来分析——一段在交 $X_0$ 内的道路，
一段在 $X_1$ 内但在 $X_2$ 外的道路，以及一段在 $X_2$ 内但在 $X_1$ 外的道路。
通常包含映射诱导的映射不是单射，因此更精确的定理版本使用群的推出（push-out）。
在群胚范畴中对非连通空间有推广（正如非交换几何中的情形），还有「高维」推广
（例如使用 2-群等；更多细节可参看 J.P. May 的代数拓扑讲义）。

接下来我们回顾同调代数中所谓的\emph{5-引理}。事实上它在任何阿贝尔范畴中都成立
（如阿贝尔群范畴，或给定域上的向量空间范畴）：

% [原文 p.128]
\para{5-引理.} 设在一个阿贝尔范畴中给定如下交换图（梯子）：
\[
  \begin{tikzcd}
    A \arrow[r,"f"] \arrow[d,"l"'] & B \arrow[r,"g"] \arrow[d,"m"'] & C \arrow[r,"h"] \arrow[d,"n"'] & D \arrow[r,"j"] \arrow[d,"p"'] & E \arrow[d,"q"] \\
    A' \arrow[r,"r"'] & B' \arrow[r,"s"'] & C' \arrow[r,"t"'] & D' \arrow[r,"u"'] & E'
  \end{tikzcd}
\]
% [Tutor Note: p.128 原文把 "ladder"（梯子）拼作 "(lader)"，少一个 d；译文按其本意译作「梯子」，并在此登记原文拼写。]
则若两行都是正合列，$m$ 与 $p$ 是同构，$l$ 是满态射且 $q$ 是单态射，则 $n$ 也是同构。

证明的方法是「图追踪」。

5-引理是两个 4-引理的组合，其中一个与另一个对偶。一个有趣的特殊情形是所谓的\emph{短 5-引理}，
其中两行是短正合列；此时有 $A = A' = E = E' = 0$，因此只有三个竖映射 $m, n, p$，
而短 5-引理断言：若 $m, p$ 是同构，则 $n$ 也是。

同调代数中一个密切相关的陈述是\emph{蛇形引理}：

\para{蛇形引理.} 在一个阿贝尔范畴中（例如给定域上的向量空间范畴），考虑如下交换图：
\[
  \begin{tikzcd}
    0 \arrow[r] & A \arrow[r,"f"] \arrow[d,"a"'] & B \arrow[r,"g"] \arrow[d,"b"'] & C \arrow[r] \arrow[d,"c"] & 0 \\
    0 \arrow[r] & A' \arrow[r,"{f'}"'] & B' \arrow[r,"{g'}"'] & C' \arrow[r] & 0
  \end{tikzcd}
\]
其中两行正合且 $0$ 是零对象。则存在一个把映射 $a, b, c$ 的核与余核联系起来的正合列：
\[
  0 \to \operatorname{Ker} a \to \operatorname{Ker} b \to \operatorname{Ker} c
  \xrightarrow{d} \coker a \to \coker b \to \coker c \to 0.
\]

% [原文 p.129]
此外，若 $f$ 是单态射，则映射 $\operatorname{Ker} a \to \operatorname{Ker} b$ 也是单态射；
若 $g'$ 是满态射，则映射 $\coker b \to \coker c$ 也是满态射。

接下来我们回顾\emph{万有系数定理}，它给出某个拓扑空间 $X$ 的整系数同调 $H_i(X; \Z)$
与以任意阿贝尔群 $A$ 为系数的相应同调 $H_i(X; A)$ 之间的关系：

\para{万有系数定理.} 沿用上面的记号，考虑张量积 $H_i(X; \Z) \otimes A$。
则存在群同态 $i : H_i(X; \Z) \otimes A \to H_i(X; A)$，它是单射，其余核是
$\operatorname{Tor}(H_{i-1}(X; A), A)$。

换言之，存在自然的短正合列
\[
  0 \to H_i(X; \Z) \otimes A \to H_i(X; A) \to \operatorname{Tor}(H_{i-1}(X; A), A) \to 0.
\]
该列分裂（但不是以自然的方式分裂），而挠群可以看作 $i$ 成为同构的障碍。
上同调有相应的（对偶）定理。

\emph{Hurewicz 定理}是代数拓扑的一个中心结果，它推广了 Poincaré 关于同伦与同调之间关系的定理。
其中一个版本如下：

\para{Hurewicz 定理.} 对任何拓扑空间 $X$ 与自然数 $k$，存在群同态
\[
  h_* : \pi_k(X) \to H_k(X; \Z)
\]
它称为\emph{Hurewicz 同态}。有相对版本，还有三元组（triadic）版本。

对 $k = 1$，Hurewicz 同态是典范的交换化映射
\[
  h_* : \pi_1(X) \to \pi_1(X)/[\pi_1(X), \pi_1(X)].
\]
特别地，第一（整系数）同调群同构于基本群的交换化。此外，若 $X$ 是 $(n-1)$-连通的，则%
% [原文 p.130]
Hurewicz 同态是同构，$\forall k \leq n$（这意味着：若空间道路连通且基本群完美，
则第一同调群为零）。

\emph{Whitehead 定理}为代数拓扑中使用 $CW$-复形提供了依据，因为它断言：
若两个拓扑空间之间的连续映射诱导了所有相应同伦群之间的同构（换言之，它是弱同伦等价），
则该映射事实上是同伦等价，当且仅当这两个空间是连通的 $CW$-复形：

\para{Whitehead 定理.} 设 $X, Y$ 是两个 $CW$-复形，其基点分别为 $x, y$，
并设 $f : X \to Y$ 是连续映射，满足 $f(x) = y$。对任何正整数 $n \neq 0$，考虑诱导同态
\[
  f_* : \pi_n(X, x) \to \pi_n(Y, y).
\]
若 $f_*$ 对所有 $n$ 都是同构，则称 $f$ 是弱同伦等价。于是若 $f$ 是弱同伦等价，
则它也是同伦等价。

Whitehead 定理对非 $CW$-复形空间的推广，由拓扑学的另一分支——所谓\emph{形状理论}——来研究。
此外，\emph{Quillen} 证明了在任何\emph{模型范畴}中，\emph{纤维性}与\emph{余纤维性}对象之间的
弱同伦等价事实上是同伦等价。

\emph{Eilenberg--MacLane 空间}是同伦理论的积木：

\para{Eilenberg--MacLane 空间的定义.} 设 $G$ 是群，$n$ 是非零自然数。
连通拓扑空间 $X$ 称为 $K(G, n)$ 型的\emph{Eilenberg--MacLane 空间}，
若第 $n$ 个同伦群 $\pi_n(X)$ 同构于 $G$，且 $X$ 的所有其他同伦群为零。

对 $n > 1$，$G$ 必须是交换群；因此作为 $CW$-复形，Eilenberg--MacLane 空间存在，
且在弱同伦意义下唯一。

% [原文 p.131]
例子：
\begin{itemize}
\item $K(\Z, 1) = S^1$。
\item $K(\Z, 2) = \C P^{\infty}$。
\item $K(\Z_2, 1) = \R P^{\infty}$。
\item $K(G, 1) = \bigwedge_{i=1}^{k} S^1$，其中 $G$ 是 $k$ 个生成元的自由群。
\end{itemize}

每个空间 $K(G, n)$ 都可以构造为一个 $CW$-复形：从若干因子 $S^n$ 的 smash 积出发
（$G$ 的每个生成元对应一个因子），然后添加更高维的胞腔（数目可能无限）以消去剩余的同伦。

Eilenberg--MacLane 空间的基本性质是：它们给出以 $G$ 为系数的同调的表示。
对任何阿贝尔群 $G$ 与 $CW$-复形 $X$，映射 $f : X \to K(G, n)$ 的同伦类之集 $[X, K(G, n)]$
与集合 $H^n(X; G)$ 自然地一一对应（后者是 $X$ 的以 $G$ 为系数的第 $n$ 个奇异上同调群）。

\para{Milnor--Moore 定理.} 设 $X$ 是 1-连通空间，环路空间为 $\Omega X$。
则 Hurewicz 同态诱导 Hopf 代数同构
\[
  U(\pi_*(\Omega X) \otimes \Q) \to H_*(\Omega X; \Q)
\]
其中 $U$ 表示泛包络代数，而 $\pi_*(\Omega X) \otimes \Q$ 的李括号由 Samelson 积给出。

\para{群完备化定理（D. McDuff, G. Segal）.} 设 $M$ 是拓扑幺半群，$BM$ 是它的分类空间。
设 $M \to \Omega BM$ 是典范映射。则映射
\[
  H_*(M) \to H_*(\Omega BM)
\]
诱导出同构
\[
  H_*(M)[\pi_0(M)^{-1}] \to H_*(\Omega BM).
\]

% [原文 p.132]
\emph{闭图像定理}是泛函分析中的一个重要结果：

\para{闭图像定理.} 设 $T : X \to Y$ 是两个 Banach 空间之间的线性算子，且在整个 $X$ 上有定义。
$T$ 的图像定义为
\[
  \{(x, y) \in X \times Y : Tx = y\}.
\]
则 $T$ 连续当且仅当它是闭的，即当且仅当它的图像是笛卡尔积 $X \times Y$ 的闭子集
（赋予乘积拓扑）。

回忆一下，分析（实分析与复分析）中的\emph{Weierstrass 定理}（1885）断言：
定义在闭区间上的每个连续函数都可以被多项式函数一致逼近。

1937 年 \emph{M.H.Stone} 从两个方向推广了上述结果：他把闭区间换成任意紧 Hausdorff 空间 $X$
（这还可以进一步推广到非紧的 \emph{Tychonoff} 空间），并把多项式函数换成 $C(X)$
（$X$ 上连续复函数代数）的一般子代数。更具体地说，我们有：

\para{Stone--Weierstrass 定理.} 设 $X$ 是紧 Hausdorff 空间，$S \subset C(X)$ 是一个
\emph{分离点}的子集（即对任何一对点 $x, y \in X$，存在函数 $f \in S$ 使 $f(x) \neq f(y)$）。
则由 $S$ 生成的复单位 $\ast$-代数在 $C(X)$ 中稠密。

注意该定理还有实版本，以及若干其他推广（涉及局部紧 Hausdorff 空间等）。

我们用 \emph{Gelfand 定理}来结束本节，它是发展非交换几何的一个重要动机：

\para{Gelfand 定理.} \emph{如下两个范畴是等价的：}
\begin{enumerate}[label=\arabic*.]
\item 对象为有单位的交换 $C^*$-代数、态射为保 $\ast$ 同态的范畴。
\item 对象为紧 Hausdorff 空间、态射为它们之间的%
% [原文 p.133]
同胚的范畴。
\end{enumerate}

第一个函子记作 $C(-)$，它把每个紧 Hausdorff 空间 $X$ 对应到 $X$ 上连续复函数的有单位交换
$C^*$-代数 $C(X)$；第二个函子记作 $\operatorname{Spec}(-)$，它把每个交换的有单位 $C^*$-代数 $A$
对应到紧 Hausdorff 空间 $\operatorname{Spec}(A)$，即 $A$ 的谱。

上述等价可以推广到局部紧 Hausdorff 空间与非单位的交换 $C^*$-代数之间。

伟大的法国数学家 Alain Connes 发展非交换几何的出发点之一，正是在考虑非交换代数时
上述等价的「拓扑推广」。

% [Tutor Note: p.127 Van Kampen 定理的记号核对：$i_k$ 由包含 $X_0\into X_k$ 诱导，$j_k$ 由包含 $X_k\into X$ 诱导（$k=1,2$）；核心里那个带下标的星号是自由积的融合记号，即 $G_1\star_{G_0}G_2$（$G_0$ 的两个像被等同）。原文随后自己解释了「通常包含诱导的映射不是单射，因此精确版本用 push-out」——这正是从「融合自由积」升级到「群范畴中的推出」的原因：只有当 $i_1,i_2$ 都是单射时，融合自由积才是范畴论意义上的推出。]
% [Tutor Note: p.127 首段 "background meterial" 拼写有误，应为 "material"（背景材料）；已按 500 dpi 裁图复核页图确为 meterial，译文按其意译出并在回执中登记。]
% [Tutor Note: p.128 5-引理的假设只需「$l$ 满、$q$ 单、$m,p$ 同构」即可推出 $n$ 同构——注意 $n$ 本身没有假设。核对方法：从 $C'$ 取元逆推。短 5-引理是 $A=A'=E=E'=0$ 的特例，此时梯子退化为 $0\to B\to C\to 0$ 型的两个短正合列，只剩三个竖映射 $m:B\to B'$、$n:C\to C'$、$p:D\to D'$（原文写法如此），断言仍是「两端同构则中间同构」。]
% [Tutor Note: p.128 蛇形引理的连接同态 $d:\operatorname{Ker} c\to\coker a$ 由「先在 $B$ 中提升、再用 $b$ 推到 $B'$、最后在 $\coker a$ 中取类」得到，其良定义性依赖两行正合；图上 $f,g$ 在上行、$f',g'$ 在下行（原文把 $f'$ 的撇排在字母右上，排版上易看成 $f^\prime$）。末尾两条附加断言也须核对：$f$ 单 $\Rightarrow\ \operatorname{Ker}a\to\operatorname{Ker}b$ 单；$g'$ 满 $\Rightarrow\ \operatorname{coker}b\to\operatorname{coker}c$ 满。]
% [Tutor Note: p.129 万有系数定理里原文把 Tor 项写成 $\operatorname{Tor}(H_{i-1}(X;A),A)$，而标准陈述是 $\operatorname{Tor}(H_{i-1}(X;\Z),A)$（系数应取整系数同调）。已用 500 dpi 裁图复核页图，原文确实两处都写 $H_{i-1}(X;A)$，疑为笔误；译文照抄原文，此处给正确形式供读者对照。]
% [Tutor Note: p.130 原文把 "CW-complex" 一律排成「$CW-\!complex$」（用减号而非连字符，且整体斜体），译文统一写作 $CW$-复形，不再保留孤立的减号。另：Whitehead 定理陈述中「positive integer $n\neq0$」是赘语（正整数本就不为 0），照译；开头那句「……则是同伦等价，当且仅当空间是连通 $CW$-复形」的条件从句位置含糊，严格的说法是「对连通的 $CW$-复形，弱同伦等价蕴含同伦等价」，译文照原文语序。]
% [Tutor Note: p.130--131 Eilenberg--MacLane 空间的例子核对：$K(\Z,1)=S^1$、$K(\Z,2)=\C P^{\infty}$、$K(\Z_2,1)=\R P^{\infty}$ 分别因为 $\pi_1(S^1)=\Z$、$\pi_2(\C P^\infty)=\Z$ 且其余同伦群为零、$\pi_1(\R P^\infty)=\Z_2$。第四个例子原文用 $\bigwedge_{i=1}^{k}S^1$（大楔积记号），实际含义是 $k$ 个圆周的一点并（wedge），只有 $G$ 为 $k$ 元自由群时才是 $K(G,1)$；$n>1$ 时 $G$ 必须交换这一点原文已说明，而 $n=1$ 时 $G$ 可任意。]
% [Tutor Note: p.131 "to kill of the remaining homotopy" 中 "kill of" 应为 "kill off"（消去），已按 500 dpi 裁图复核；译文按「消去」处理。同段「它们给出以 $G$ 为系数的同调的表示」与紧随其后的 $H^n(X;G)$（奇异上同调）不一致：Eilenberg--MacLane 空间表示的是上同调 $H^n(-;G)\cong[X,K(G,n)]$，原文此处 "homology" 疑为 "cohomology" 之误，译文照抄并在此说明。]
% [Tutor Note: p.132 Stone--Weierstrass 定理中 "seperates points" 拼写有误，应为 "separates points"（分离点）；已按 500 dpi 裁图复核。该条件等价于「$S$ 中函数能区分任意两个不同点」，是 $S$ 生成的代数在 $C(X)$ 中稠密的关键假设之一（另需 $S$ 对共轭封闭，原文未提，属省略）。]
% [Tutor Note: p.132 Gelfand 定理第 2 条的态射原文写作「arrows the homeomorphisms among them」，但按 Gelfand 对偶，紧 Hausdorff 空间范畴的态射应是连续映射（$f:X\to Y$ 对应 $f^*:C(Y)\to C(X)$）；同胚只对应 $\ast$-同构，若只取同胚则得到的是「范畴的等价于 $\ast$-代数范畴的骨架」这类弱陈述。判断依据：p.133 立刻用反变函子 $C(-)$ 与 $\operatorname{Spec}(-)$ 说明等价，反变函子把任意连续映射送到 $\ast$-同态，而不是只把同胚送到同构。译文照抄原文并在回执登记。]
% [Tutor Note: p.133 $C(-)$ 与 $\operatorname{Spec}(-)$ 都是反变函子：$X\mapsto C(X)$ 把连续映射 $f:X\to Y$ 送到 $f^*:C(Y)\to C(X)$，$A\mapsto\operatorname{Spec}(A)$ 把 $\ast$-同态送到其诱导的谱映射；两者复合与恒等函子自然同构，这才是「两个范畴等价」的确切含义。原文说该等价可推广到「局部紧 Hausdorff 空间 ↔ 非单位交换 $C^*$-代数」，推广时要用 $C_0(X)$（见 p.76 旁注 1）。]
